\documentclass{article}
\usepackage{booktabs}
\usepackage{amsmath}
\usepackage{geometry}
\usepackage{tabularray}
\usepackage{float}
\usepackage{graphicx}
\usepackage{codehigh}
\usepackage[normalem]{ulem}
\UseTblrLibrary{booktabs}
\UseTblrLibrary{siunitx}
\newcommand{\tinytableTabularrayUnderline}[1]{\underline{#1}}
\newcommand{\tinytableTabularrayStrikeout}[1]{\sout{#1}}
\NewTableCommand{\tinytableDefineColor}[3]{\definecolor{#1}{#2}{#3}}
\geometry{margin=1in}

\title{\textbf{Problem Set 8: } \\
Matrix Algebra and Optimization}

\date{\today}

\begin{document}
\maketitle

\section*{Question 5: Closed-Form OLS}
The closed-form OLS estimator $\hat{\beta}^{OLS} = (X'X)^{-1}X'Y$ yields estimates 
very close to the true $\beta$ vector. With $N = 100{,}000$ observations, the law of 
large numbers ensures the estimates converge tightly to the true values. Specifically, 
the maximum deviation from the truth across all ten coefficients is less than $0.01$.

\section*{Question 6: Gradient Descent}
Using a learning rate of $0.0000003$, gradient descent converged at iteration 274, 
recovering estimates that match the closed-form solution to at least four decimal places. 
The small step size ensures numerical stability throughout the iterations. This confirms 
that the OLS objective function is convex and gradient descent reliably finds the global 
minimum.

\section*{Question 7: L-BFGS and Nelder-Mead}
Both the L-BFGS and Nelder-Mead algorithms recover coefficient estimates essentially 
identical to the closed-form OLS solution. The two methods do not materially differ 
from each other, confirming that the objective function is well-behaved and the 
optimum is unique. L-BFGS uses gradient information while Nelder-Mead is 
derivative-free, yet both converge to the same answer.

\section*{Question 8: MLE via L-BFGS}
The MLE estimates of $\beta$ match the OLS estimates exactly, as expected in a linear 
model with normally distributed errors where OLS and MLE coincide. The estimated 
$\hat{\sigma} = 0.4999$, which is extremely close to the true value of $\sigma = 0.5$ 
used to generate the data. This provides further validation that the data generating 
process was correctly implemented.

\section*{Question 9: OLS via \texttt{lm()} and Comparison to Ground Truth}
Table 1 presents the OLS estimates obtained using \texttt{lm(Y \textasciitilde\ X - 1)}, 
where the \texttt{-1} suppresses R's automatic intercept since the first column of $X$ 
already contains a vector of ones. All ten coefficient estimates are extremely close 
to the true $\beta$ values:
\[
\beta = [1.5,\ -1,\ -0.25,\ 0.75,\ 3.5,\ -2,\ 0.5,\ 1,\ 1.25,\ 2]'
\]
The maximum deviation from the truth across all ten coefficients is less than $0.01$, 
which reflects the precision gained from a large sample of $N = 100{,}000$ observations. 
This confirms that OLS is unbiased and consistent under the classical linear model 
assumptions. All methods --- closed-form, gradient descent, L-BFGS, Nelder-Mead, 
and MLE --- recover virtually identical estimates, further validating the implementation.

\input{PS8_table_Jhinuk.tex}

\end{document}